% ============================================================
\section{Dữ liệu và knowledge base}
\label{sec:data}
% ============================================================

\subsection{Knowledge base do giảng viên cung cấp}

KB gồm sáu file văn bản trong \code{data/kb/} (\cref{tab:kb}). Đây là nguồn dữ kiện
duy nhất của câu trả lời; nhóm không bổ sung hay sửa dữ kiện KB.

\begin{table}[htbp]
  \centering
  \small
  \begin{tabularx}{\linewidth}{|P{3.4cm}|L|P{3.6cm}|}
    \hline
    \textbf{File} & \textbf{Nội dung} & \textbf{Collection sau khi nạp} \\ \hline
    \code{course_info.txt} & Mã, tên môn, học kỳ, tín chỉ, mô tả, LO, cấu trúc đánh
      giá, hình thức thi giữa kỳ và cuối kỳ & \code{course} \\ \hline
    \code{schedule.txt} & Lịch 15 tuần: chương, chủ đề, hoạt động và mốc bài tập lớn
      của từng tuần & \code{weeks} \\ \hline
    \code{topics.txt} & 12 chương: phần kiến thức, tóm tắt, chủ đề con, từ khóa &
      \code{topics}, \code{chapter_weeks} \\ \hline
    \code{assignments.txt} & Bài tập lớn BTL01: mục tiêu, yêu cầu tối thiểu của Phần
      I--IV, phần mở rộng, intent và output gợi ý & \code{assignment} \\ \hline
    \code{resources.txt} & 8 tài liệu tham khảo: tên, tác giả, năm, cách sử dụng &
      \code{resources} \\ \hline
    \code{regulations.txt} & Quy định: chuyên cần, liêm chính, sử dụng AI, thư viện
      ngoài, LLM/RAG, câu hỏi ngoài phạm vi & \code{regulations} \\ \hline
  \end{tabularx}
  \caption{Sáu file KB và collection tương ứng}
  \label{tab:kb}
\end{table}

Bộ nạp KB (\code{kb.py}) đọc từng file thành collection có kiểu, chuẩn hóa định danh
(ví dụ chương 3 có ID \code{CH03}, tuần 3 có ID \code{WEEK_03}, \code{chapter_number}
là số nguyên), lập chỉ mục \code{chapter_weeks} từ chương sang các tuần học, và kiểm
tra tính nhất quán giữa lịch học với danh sách chương. Nhóm lưu một bản golden của KB
đã nạp để phát hiện mọi thay đổi về dữ kiện hoặc kiểu dữ liệu khi sửa bộ nạp.

\paragraph{Dữ kiện KB không có.} KB không chứa tên giảng viên, ngày thi cụ thể, số buổi
học mỗi tuần, danh sách output riêng cho Phần II--IV, cũng như quan hệ chi tiết giữa
tài liệu và chương. Với các câu hỏi về những trường này, hệ thống vẫn phân tích được
câu và xác định intent, nhưng truy vấn trả về rỗng và câu trả lời là thông báo không
tìm thấy. Nhóm không điền các trường này bằng suy đoán.

\subsection{Dữ liệu hỗ trợ và dữ liệu kiểm thử}

Thư mục \code{data/scaffolding/} chứa phần dữ liệu khung của đề và dữ liệu nhóm bổ sung:
\begin{itemize}
  \item \code{entities.txt}: lược đồ thực thể và bảng alias (ví dụ ``NLP'', ``môn này''
        $\rightarrow$ \code{CO3085}; ``BTL'' $\rightarrow$ \code{BTL01}).
  \item \code{sample_queries.txt}: 42 câu hỏi mẫu của đề, kèm intent, entity, trạng
        thái truy vấn và các fact bắt buộc trong câu trả lời.
  \item \code{dialogues.txt}: 6 hội thoại mẫu của đề dùng cho Phần III.
  \item \code{challenge_queries.txt}: 22 câu thử thách \textbf{do nhóm tự biên soạn}
        để phân tích lỗi và giới hạn; header của file ghi rõ nguồn này.
\end{itemize}

File \code{input/sentences.txt} chứa 73 câu đầu vào cho lệnh \code{run.py all}: các câu
mẫu của đề và các cách diễn đạt khác do nhóm viết trong giai đoạn 1. Dòng trống và
dòng chú thích bắt đầu bằng \texttt{\#} được bỏ qua. Dữ liệu bổ sung chỉ dùng để kiểm
thử, không thay đổi bản chất bài toán và không tạo câu trả lời ngoài KB \cite{spec}.
